\originalchapter{谱函数与Fermi液体}
\label{ch:fermi}
\source{8.512第23--25讲。}
\originalsection{单粒子谱的物理含义}
推迟电子Green函数定义$G^R(k,t)=-(\ii/\hbar)\theta(t)\avg{\{c_k(t),c_k^\dagger(0)\}}$。Fourier变换后$G^R(k,\omega)=[\hbar\omega-\epsilon_k+\mu-\Sigma^R(k,\omega)]^{-1}$，谱函数$A(k,E)=-\pi^{-1}\operatorname{Im}G^R(k,E/\hbar)$满足$\int A(k,E)\dd E=1$。占据数$n_k=\int f(E)A(k,E)\dd E$。光电子发射主要测$fA$并乘矩阵元；隧穿测动量积分，二者并非同一观测量。

若自能在费米面附近可展成$\Sigma^R\simeq\Sigma_0+(1-Z^{-1})E+\delta k\,\partial_k\Sigma-\ii C[E^2+(\pi\kb T)^2]$，则
\begin{equation}
G^R\simeq\frac Z{E-v_F^*\hbar(k-k_F)+\ii\Gamma(E,T)}+G_{
m incoh}.
\end{equation}
$Z=[1-\partial_E\Sigma']^{-1}$是单粒子极点权重，不是总粒子数比例。零温$n_k$在费米面有跳跃$Z$，其余权重在非相干连续谱中。衰减率$\Gamma\propto E^2$使$\Gamma/E\to0$，低能激发才有可分辨的准粒子意义。若自能非解析导致$Z=0$或$\Gamma/E$不消失，不能继续使用此极点展开。
\originalsection{能量泛函与Landau参数}
Fermi液体的状态以准粒子分布变化$\delta n_{k\sigma}$标记，能量展开
\begin{equation}
\delta E=\sum_{k\sigma}\epsilon^*_{k}\delta n_{k\sigma}+\frac1{2V}\sum_{kk'\sigma\sigma'}f_{\sigma\sigma'}(k,k')\delta n_{k\sigma}\delta n_{k'\sigma'}.
\end{equation}
$f$是改变其他准粒子分布时一个准粒子能量的变化，不等同于裸库仑势。各向同性、自旋旋转对称时分成$s,a$通道，再在夹角的Legendre多项式上展开，定义$F_l^{s,a}=D^*(E_F)f_l^{s,a}$，本文$D^*$为双自旋态密度。均匀密度与自旋响应为
\begin{equation}
\frac{\partial n}{\partial\mu}=\frac{D^*(E_F)}{1+F_0^s},\qquad \chi_s=\frac{\mu_0\mu_B^2D^*(E_F)}{1+F_0^a},\qquad\frac{C_V}{V}=\frac{\pi^2}3D^*(E_F)\kb^2T.
\end{equation}
比热只测热态数，压缩率和磁化率还含分布自洽调整，所以质量增强不足以预测它们。

在Galilean不变连续系统，均匀整体boost的电流必须由裸质量决定，准粒子的附带回流给$m^*/m=1+F_1^s/3$。在晶格上Galilean不变性不存在，这条质量关系一般失效。稳定性要求各角动量通道$1+F_l^{s,a}/(2l+1)>0$；某个系数变负意味着相应费米面形变可降低能量，即Pomeranchuk不稳定。
\originalsection{碰撞、零声与Kondo低能态}
线性化无碰撞动力学$\partial_t\delta n+\vct v_F\cdot\nabla_r[\delta n-f'(\epsilon)\delta\epsilon]=0$，相互作用使密度畸变自洽。只保留$F_0^s$的三维纵模令$s=\omega/(qv_F)$，角积分给
\begin{equation}
1+F_0^s\left[1-\frac s2\log\frac{s+1}{s-1}\right]=0.
\end{equation}
排斥$F_0^s>0$有$s>1$解，模在粒子--空穴连续谱外，是零声。碰撞足够频繁$\omega\tau\ll1$时局域平衡建立，得到普通第一声；两个声不是相同极限。粒子--空穴相空间在低温宽度$T$内提供两个独立能量积分，所以准粒子碰撞率$\propto T^2$，有限能量对应$E^2$。

Kondo单重态附近，杂质散射相移可展开$\delta_\sigma(E,\delta n)=\pi/2+\alpha E/(\kb T_K)-\phi\,\delta n_{-\sigma}/[D_\uparrow\kb T_K]+\cdots$。能量项与残余相互作用项共同决定磁化率和比热，局域液体的Wilson比在标准完全屏蔽极限趋2。该结果需要强耦合固定点分析，前章一回路对数流只给$T_K$尺度，不能单独证明精确低温比值。
